% ==================================================
% Chapter BACKGROUND
% ==================================================

\section{背景}
某“声学监测仪”长期部署在近岸浅海区域，采用24V铅酸电池组作为核心供电单元。由于海洋环境复杂，设备投放和回收周期长，电池需要具备稳定的快充能力和可靠的安全保障：一方面，为了缩短设备回收后的维护时间，需要在保证电池寿命的前提下实现最大充电速率，安全充电电流需严格控制在5A；另一方面，浅海环境下电压波动频繁，电池过充、过流会直接导致传感器失灵，甚至引发设备短路损坏，导致监测设备在水下停机的情况。

为满足这一需求，学院创新实验室XX决定以项目实战为导向，开展硬件核心成员招新考核。

本次考核的铅酸电池充电电路设计任务，正是源于“浅海便携式声学监测系统”的实际需求：要求设计的充电管理装置，不仅要满足实验室现有设备的限制（输入直流电压不超过48V），还要适配水下设备的严苛标准——低内阻、高精度电流检测、快速响应的过流保护，以及充满自动停止等功能。


